% ===== PRÉAMBULE CANONIQUE — à coller VERBATIM en tête de CHAQUE .tex =====
\documentclass[11pt,a4paper]{article}
\usepackage[utf8]{inputenc}\usepackage[T1]{fontenc}\usepackage{lmodern}
\usepackage[french]{babel}
\usepackage{amsmath,amssymb,mathtools}\usepackage{icomma}
\usepackage[version=4]{mhchem}          % \ce{...} : équations & formules
\usepackage{siunitx}\sisetup{locale=FR} % \qty{}{}, \si{}, \ang{}
\usepackage{chemfig}                    % \chemfig{...} : structures développées
\usepackage{xcolor}\usepackage{colortbl}
\usepackage{tikz}\usepackage{pgfplots}\pgfplotsset{compat=1.18}
\usepackage[most]{tcolorbox}\usepackage{enumitem}\usepackage{array,booktabs}
\usepackage[a4paper,margin=2.2cm]{geometry}
\usetikzlibrary{arrows.meta,calc,angles,quotes,positioning,patterns,backgrounds,shapes.geometric}
\definecolor{primary}{RGB}{222,49,99}\definecolor{primarydark}{RGB}{150,22,55}
\definecolor{defcol}{RGB}{0,114,178}\definecolor{methcol}{RGB}{52,92,148}
\definecolor{excol}{RGB}{0,135,90}\definecolor{attncol}{RGB}{200,40,55}
\tcbset{boxbase/.style={enhanced,breakable,boxrule=0.9pt,arc=2.5pt,
  left=3mm,right=3mm,top=2.3mm,bottom=2.3mm,fonttitle=\bfseries,coltitle=white}}
% --- boîtes TUTORIEL (noms EXACTS, ne pas renommer) ---
\newtcolorbox{cle}[1][]{boxbase,colback=defcol!6,colframe=defcol,
  title={\textsc{À retenir}\ifx\relax#1\relax\relax\else\textmd{ : \itshape #1}\fi}}
\newtcolorbox{methode}[1][]{boxbase,colback=methcol!7,colframe=methcol,sharp corners,
  title={\textsc{Méthode}\ifx\relax#1\relax\relax\else\textmd{ --- \itshape #1}\fi}}
\newtcolorbox{exemplebox}[1][]{boxbase,colback=excol!5,colframe=excol,
  title={\textsc{Exemple}\ifx\relax#1\relax\relax\else\textmd{ : \itshape #1}\fi}}
\newtcolorbox{piege}[1][]{boxbase,colback=attncol!6,colframe=attncol,
  title={\textsc{Piège}\ifx\relax#1\relax\relax\else\textmd{ : \itshape #1}\fi}}
\newtcolorbox{remarque}[1][]{boxbase,colback=black!4,colframe=black!45,
  title={\textsc{Remarque}\ifx\relax#1\relax\relax\else\textmd{ : \itshape #1}\fi}}
% --- boîtes EXERCICE / CORRIGÉ (noms EXACTS, auto-comptées) ---
\newtcolorbox[auto counter]{exo}[1][]{boxbase,colback=primary!4,colframe=primary,
  title={\textsc{Exercice}~\thetcbcounter\ifx\relax#1\relax\relax\else\textmd{ --- \itshape #1}\fi}}
\newtcolorbox[auto counter]{corrige}[1][]{boxbase,colback=excol!4,colframe=excol,
  title={\textsc{Corrigé}~\thetcbcounter\ifx\relax#1\relax\relax\else\textmd{ --- \itshape #1}\fi}}
\newcommand{\R}{\mathbb{R}}\newcommand{\N}{\mathbb{N}}\newcommand{\Z}{\mathbb{Z}}\newcommand{\ds}{\displaystyle}
% (un \usepackage{fancyhdr}+en-tête est possible ; il est ignoré côté web)
% =========================================================================
\begin{document}
% THEME_TITLE: Incertitude, justesse et fidélité
\sisetup{per-mode=symbol}

Les chiffres significatifs disent la \emph{précision d'écriture} d'un nombre. Mais une mesure peut être
précise sans être correcte : il faut deux notions de plus, et une façon de \emph{chiffrer} le doute.

\begin{cle}[justesse et fidélité]
\begin{itemize}[leftmargin=1.6em,topsep=2pt,itemsep=2pt]
  \item \textbf{Justesse} : proximité de la mesure à la \emph{vraie} valeur (mesure non biaisée).
  \item \textbf{Fidélité} : proximité des mesures \emph{entre elles} quand on répète (mesure reproductible).
\end{itemize}
\end{cle}

\begin{center}
\begin{tikzpicture}[
   ring/.pic={
     \fill[attncol!12] (0,0) circle (1.2); \fill[white] (0,0) circle (0.8);
     \fill[attncol!12] (0,0) circle (0.4);
     \draw[black!55] (0,0) circle (1.2); \draw[black!55] (0,0) circle (0.8);
     \draw[black!55] (0,0) circle (0.4); \fill[black!65] (0,0) circle (0.045);}]
  \begin{scope}
     \pic{ring};
     \fill[primary](0,0)circle(0.075);\fill[primary](0.13,0.09)circle(0.075);
     \fill[primary](-0.09,0.12)circle(0.075);\fill[primary](0.07,-0.11)circle(0.075);
     \node[below,align=center,font=\scriptsize] at (0,-1.55){juste\\ \& fidèle};
  \end{scope}
  \begin{scope}[xshift=3.4cm]
     \pic{ring};
     \fill[primary](0.75,0.62)circle(0.075);\fill[primary](0.88,0.72)circle(0.075);
     \fill[primary](0.64,0.72)circle(0.075);\fill[primary](0.80,0.52)circle(0.075);
     \node[below,align=center,font=\scriptsize] at (0,-1.55){fidèle,\\ non juste};
  \end{scope}
  \begin{scope}[xshift=6.8cm]
     \pic{ring};
     \fill[primary](-0.72,0.22)circle(0.075);\fill[primary](0.62,-0.45)circle(0.075);
     \fill[primary](0.16,0.75)circle(0.075);\fill[primary](-0.30,-0.70)circle(0.075);
     \node[below,align=center,font=\scriptsize] at (0,-1.55){juste,\\ non fidèle};
  \end{scope}
  \begin{scope}[xshift=10.2cm]
     \pic{ring};
     \fill[primary](0.70,0.78)circle(0.075);\fill[primary](1.00,0.42)circle(0.075);
     \fill[primary](-0.82,-0.30)circle(0.075);\fill[primary](0.30,-0.92)circle(0.075);
     \node[below,align=center,font=\scriptsize] at (0,-1.55){ni juste\\ ni fidèle};
  \end{scope}
\end{tikzpicture}
\end{center}

\begin{cle}[incertitude absolue]
On écrit une mesure sous la forme $x=x_{\text{mes}}\pm\Delta x$, où $\Delta x$ (même unité que $x$) est
l'«~épaisseur de doute~». Par convention, le \emph{dernier chiffre significatif écrit} fixe l'ordre de
$\Delta x$.
\end{cle}

\begin{exemplebox}[de l'écriture à l'incertitude]
$m=\qty{13,02}{\gram}$ sous-entend $\Delta m=\qty{0,01}{\gram}$, donc $m=(13,02\pm0,01)\ \si{\gram}$ : la
vraie masse est comprise entre $\qty{13,01}{\gram}$ et $\qty{13,03}{\gram}$.
\end{exemplebox}

\begin{cle}[incertitude relative]
\[ \frac{\Delta x}{x}\quad\text{(sans unité, souvent en \%)}. \]
Elle mesure le \emph{poids relatif} du doute et permet de comparer la qualité de mesures de grandeurs
différentes.
\end{cle}

\begin{exemplebox}[lien chiffres significatifs $\leftrightarrow$ incertitude relative]
Le nombre de CS et l'incertitude relative vont de pair. Ordres de grandeur typiques :
\begin{center}\renewcommand{\arraystretch}{1.15}
\begin{tabular}{c c c c}
\toprule
\textbf{CS} & $2$ & $3$ & $4$\\
\midrule
\textbf{incertitude relative} & $\sim 1\%$ & $\sim 0,1\%$ & $\sim 0,01\%$\\
\bottomrule
\end{tabular}
\end{center}
Plus il y a de CS, plus l'incertitude relative est faible : le nombre de CS \emph{est} une lecture
directe de la précision relative.
\end{exemplebox}

\begin{piege}[les CS mesurent la fidélité, pas la justesse]
Une balance déréglée peut afficher $\qty{13,02}{\gram}$ ($4$ CS, très «~précis~») tout en étant
\emph{fausse} de $\qty{2}{\gram}$. Les CS traduisent la \emph{fidélité} (finesse de lecture), jamais la
\emph{justesse} : ils ne corrigent aucun biais.
\end{piege}

\end{document}
